% Options for packages loaded elsewhere
\PassOptionsToPackage{unicode}{hyperref}
\PassOptionsToPackage{hyphens}{url}
\documentclass[
]{article}
\usepackage{xcolor}
\usepackage[margin=1in]{geometry}
\usepackage{amsmath,amssymb}
\setcounter{secnumdepth}{-\maxdimen} % remove section numbering
\usepackage{iftex}
\ifPDFTeX
  \usepackage[T1]{fontenc}
  \usepackage[utf8]{inputenc}
  \usepackage{textcomp} % provide euro and other symbols
\else % if luatex or xetex
  \usepackage{unicode-math} % this also loads fontspec
  \defaultfontfeatures{Scale=MatchLowercase}
  \defaultfontfeatures[\rmfamily]{Ligatures=TeX,Scale=1}
\fi
\usepackage{lmodern}
\ifPDFTeX\else
  % xetex/luatex font selection
\fi
% Use upquote if available, for straight quotes in verbatim environments
\IfFileExists{upquote.sty}{\usepackage{upquote}}{}
\IfFileExists{microtype.sty}{% use microtype if available
  \usepackage[]{microtype}
  \UseMicrotypeSet[protrusion]{basicmath} % disable protrusion for tt fonts
}{}
\makeatletter
\@ifundefined{KOMAClassName}{% if non-KOMA class
  \IfFileExists{parskip.sty}{%
    \usepackage{parskip}
  }{% else
    \setlength{\parindent}{0pt}
    \setlength{\parskip}{6pt plus 2pt minus 1pt}}
}{% if KOMA class
  \KOMAoptions{parskip=half}}
\makeatother
\usepackage{color}
\usepackage{fancyvrb}
\newcommand{\VerbBar}{|}
\newcommand{\VERB}{\Verb[commandchars=\\\{\}]}
\DefineVerbatimEnvironment{Highlighting}{Verbatim}{commandchars=\\\{\}}
% Add ',fontsize=\small' for more characters per line
\usepackage{framed}
\definecolor{shadecolor}{RGB}{248,248,248}
\newenvironment{Shaded}{\begin{snugshade}}{\end{snugshade}}
\newcommand{\AlertTok}[1]{\textcolor[rgb]{0.94,0.16,0.16}{#1}}
\newcommand{\AnnotationTok}[1]{\textcolor[rgb]{0.56,0.35,0.01}{\textbf{\textit{#1}}}}
\newcommand{\AttributeTok}[1]{\textcolor[rgb]{0.13,0.29,0.53}{#1}}
\newcommand{\BaseNTok}[1]{\textcolor[rgb]{0.00,0.00,0.81}{#1}}
\newcommand{\BuiltInTok}[1]{#1}
\newcommand{\CharTok}[1]{\textcolor[rgb]{0.31,0.60,0.02}{#1}}
\newcommand{\CommentTok}[1]{\textcolor[rgb]{0.56,0.35,0.01}{\textit{#1}}}
\newcommand{\CommentVarTok}[1]{\textcolor[rgb]{0.56,0.35,0.01}{\textbf{\textit{#1}}}}
\newcommand{\ConstantTok}[1]{\textcolor[rgb]{0.56,0.35,0.01}{#1}}
\newcommand{\ControlFlowTok}[1]{\textcolor[rgb]{0.13,0.29,0.53}{\textbf{#1}}}
\newcommand{\DataTypeTok}[1]{\textcolor[rgb]{0.13,0.29,0.53}{#1}}
\newcommand{\DecValTok}[1]{\textcolor[rgb]{0.00,0.00,0.81}{#1}}
\newcommand{\DocumentationTok}[1]{\textcolor[rgb]{0.56,0.35,0.01}{\textbf{\textit{#1}}}}
\newcommand{\ErrorTok}[1]{\textcolor[rgb]{0.64,0.00,0.00}{\textbf{#1}}}
\newcommand{\ExtensionTok}[1]{#1}
\newcommand{\FloatTok}[1]{\textcolor[rgb]{0.00,0.00,0.81}{#1}}
\newcommand{\FunctionTok}[1]{\textcolor[rgb]{0.13,0.29,0.53}{\textbf{#1}}}
\newcommand{\ImportTok}[1]{#1}
\newcommand{\InformationTok}[1]{\textcolor[rgb]{0.56,0.35,0.01}{\textbf{\textit{#1}}}}
\newcommand{\KeywordTok}[1]{\textcolor[rgb]{0.13,0.29,0.53}{\textbf{#1}}}
\newcommand{\NormalTok}[1]{#1}
\newcommand{\OperatorTok}[1]{\textcolor[rgb]{0.81,0.36,0.00}{\textbf{#1}}}
\newcommand{\OtherTok}[1]{\textcolor[rgb]{0.56,0.35,0.01}{#1}}
\newcommand{\PreprocessorTok}[1]{\textcolor[rgb]{0.56,0.35,0.01}{\textit{#1}}}
\newcommand{\RegionMarkerTok}[1]{#1}
\newcommand{\SpecialCharTok}[1]{\textcolor[rgb]{0.81,0.36,0.00}{\textbf{#1}}}
\newcommand{\SpecialStringTok}[1]{\textcolor[rgb]{0.31,0.60,0.02}{#1}}
\newcommand{\StringTok}[1]{\textcolor[rgb]{0.31,0.60,0.02}{#1}}
\newcommand{\VariableTok}[1]{\textcolor[rgb]{0.00,0.00,0.00}{#1}}
\newcommand{\VerbatimStringTok}[1]{\textcolor[rgb]{0.31,0.60,0.02}{#1}}
\newcommand{\WarningTok}[1]{\textcolor[rgb]{0.56,0.35,0.01}{\textbf{\textit{#1}}}}
\usepackage{longtable,booktabs,array}
\newcounter{none} % for unnumbered tables
\usepackage{calc} % for calculating minipage widths
% Correct order of tables after \paragraph or \subparagraph
\usepackage{etoolbox}
\makeatletter
\patchcmd\longtable{\par}{\if@noskipsec\mbox{}\fi\par}{}{}
\makeatother
% Allow footnotes in longtable head/foot
\IfFileExists{footnotehyper.sty}{\usepackage{footnotehyper}}{\usepackage{footnote}}
\makesavenoteenv{longtable}
\usepackage{graphicx}
\makeatletter
\newsavebox\pandoc@box
\newcommand*\pandocbounded[1]{% scales image to fit in text height/width
  \sbox\pandoc@box{#1}%
  \Gscale@div\@tempa{\textheight}{\dimexpr\ht\pandoc@box+\dp\pandoc@box\relax}%
  \Gscale@div\@tempb{\linewidth}{\wd\pandoc@box}%
  \ifdim\@tempb\p@<\@tempa\p@\let\@tempa\@tempb\fi% select the smaller of both
  \ifdim\@tempa\p@<\p@\scalebox{\@tempa}{\usebox\pandoc@box}%
  \else\usebox{\pandoc@box}%
  \fi%
}
% Set default figure placement to htbp
\def\fps@figure{htbp}
\makeatother
\setlength{\emergencystretch}{3em} % prevent overfull lines
\providecommand{\tightlist}{%
  \setlength{\itemsep}{0pt}\setlength{\parskip}{0pt}}
\usepackage{bookmark}
\IfFileExists{xurl.sty}{\usepackage{xurl}}{} % add URL line breaks if available
\urlstyle{same}
\hypersetup{
  hidelinks,
  pdfcreator={LaTeX via pandoc}}

\author{}
\date{\vspace{-2.5em}}

\begin{document}

\section{FOCEI in dMod:
Brainstorming}\label{focei-in-dmod-brainstorming}

Living document. Wir bauen das inkrementell aus, während wir die
Designentscheidungen abklopfen. Noch kein Plan, noch kein Code.

\begin{center}\rule{0.5\linewidth}{0.5pt}\end{center}

\subsection{1. Ziel}\label{ziel}

FOCEI (First-Order Conditional Estimation with Interaction) für
nichtlineare gemischte Effekte in dMod implementieren. Skopus aktuell:

\begin{itemize}
\tightlist
\item
  Nur FOCEI, keine Stufenleiter über FO/FOCE.
\item
  Allgemeine \(\Omega\) mit voller Kovarianzstruktur,
  Cholesky-parametrisiert.
\item
  Bestehende dMod-Likelihood- und Fehlermodellinfrastruktur
  (\texttt{normL2}, \texttt{err}) wiederverwenden.
\item
  Outer-Optimierung mit \texttt{trust()}.
\item
  FOCEI-Marginal-Likelihood-Approximation symbolisch ableiten,
  Modell-/Sensitivitätspipeline bleibt wie sie ist.
\item
  PEtab-Anbindung später.
\end{itemize}

\begin{center}\rule{0.5\linewidth}{0.5pt}\end{center}

\subsection{2. Theorie}\label{theorie}

\subsubsection{2.1 Modell}\label{modell}

Für \(N\) Subjekte \(i = 1, \dots, N\) mit Beobachtungen \(y_{ij}\) zu
Zeiten \(t_{ij}\): \[
y_{ij} = f(t_{ij}, \varphi_i) + \varepsilon_{ij}, \qquad \varphi_i = h(\theta, \eta_i)
\] mit - \(\theta\): feste Effekte (Populationsparameter), -
\(\eta_i \sim \mathcal{N}(0, \Omega)\): zufällige Effekte (individuelle
Abweichungen), -
\(\varepsilon_{ij} \sim \mathcal{N}\!\bigl(0, \Sigma_{ij}\bigr)\):
Residualfehler. \(\Sigma_{ij}\) ist ein \textbf{beliebiger
benutzerdefinierter Ausdruck}, kann von \(f_{ij}\), \(\theta\) und
\(\eta_i\) abhängen. Wir machen keine Annahme über die Form. Wird über
das \texttt{err}-Argument in \texttt{normL2} symbolisch kompiliert.

Die individuelle Parametrisierung \(\varphi_i = h(\theta, \eta_i)\) ist
ebenfalls beliebig (z.B. \(\varphi_i = \theta \exp(\eta_i)\)). Wird in
dMod heute schon durch \texttt{Pexpl} abgebildet.

\subsubsection{2.2 Marginale Likelihood}\label{marginale-likelihood}

Pro Subjekt: \[
L_i(\theta, \Omega) = \int p(y_i \mid \eta_i, \theta)\, p(\eta_i \mid \Omega)\, \mathrm{d}\eta_i.
\] Für nichtlineares \(f\) analytisch nicht lösbar. Verschiedene
Approximationen:

{\def\LTcaptype{none} % do not increment counter
\begin{longtable}[]{@{}
  >{\raggedright\arraybackslash}p{(\linewidth - 2\tabcolsep) * \real{0.5000}}
  >{\raggedright\arraybackslash}p{(\linewidth - 2\tabcolsep) * \real{0.5000}}@{}}
\toprule\noalign{}
\begin{minipage}[b]{\linewidth}\raggedright
Methode
\end{minipage} & \begin{minipage}[b]{\linewidth}\raggedright
Idee
\end{minipage} \\
\midrule\noalign{}
\endhead
\bottomrule\noalign{}
\endlastfoot
FO & Linearisierung um \(\eta_i = 0\) \\
FOCE & Innen \(\eta_i^* = \arg\max p(y_i\mid\eta_i)\,p(\eta_i)\) pro
Subjekt, dann Linearisierung um \(\eta_i^*\) \\
FOCEI & FOCE plus Interaktionsterm (wenn \(\Sigma\) von \(\eta\)
abhängt, z.B. proportionaler Fehler) \\
Laplace & 2. Ordnung, volle Hessische am Mode \\
SAEM & Stochastische Approximation des EM-Algorithmus,
Monte-Carlo-basiert \\
\end{longtable}
}

\subsubsection{2.3 FOCEI Zielfunktion}\label{focei-zielfunktion}

Pro Subjekt, ausgewertet bei \(\eta_i^*\): \[
-2\log L_i \;\approx\; \underbrace{\sum_j \!\left[\frac{(y_{ij} - f_{ij})^2}{\Sigma_{ij}} + \log \Sigma_{ij}\right]}_{\text{Residualbeitrag}}
\;+\; \underbrace{\eta_i^{*\top} \Omega^{-1} \eta_i^* + \log|\Omega|}_{\text{Prior auf } \eta}
\;+\; \log|H_i|
\] mit Informationsmatrix \[
H_i \;=\; G_i^\top \Sigma_i^{-1} G_i \;+\; \Omega^{-1} \;+\; \text{(Interaktion via } \partial \Sigma / \partial \eta\text{)}, \qquad
G_i = \frac{\partial f}{\partial \eta}\bigg|_{\eta_i^*}.
\]

Außen wird über \((\theta, \Omega, \Sigma)\) minimiert, innen pro
Subjekt \(\eta_i^*\) gesucht. Bilevel.

\subsubsection{2.4 Prozeduraler Ablauf, und Abgrenzung gegen
EM}\label{prozeduraler-ablauf-und-abgrenzung-gegen-em}

FOCEI ist \textbf{geschachtelt}, nicht alternierend. Pro Outer-Schritt
wird das Inner für jedes Subjekt voll konvergiert, \emph{bevor} das
Outer überhaupt einen Schritt macht.

\begin{verbatim}
Initialisiere (theta, Omega, Sigma) und eta_i = 0 fuer alle i.

repeat   // Outer-Schleife (trust)
  // Outer fragt OFV an aktueller (theta, Omega, Sigma) an
  for i = 1..N:                                      // Inner pro Subjekt
    eta_i* <- argmin_eta J(theta, eta, Omega, Sigma; data_i)
              via inneres trust(), warmstart aus letzter Outer-Iter
    H_i    <- Hessian von J an eta_i* (faellt aus trust raus)
  end
  OFV <- sum_i [ J_i(eta_i*)  +  log|H_i| ]          // FOCEI-Aggregation
  dOFV <- analytischer Gradient nach (theta, Omega, Sigma)
  // Outer verarbeitet (OFV, dOFV) und macht einen Schritt
until Outer konvergiert
\end{verbatim}

Der entscheidende Punkt: das Inner wird in jeder Outer-Iteration
\emph{vollständig} gelöst. Es gibt keinen ``burn-in dann
alternierend''-Modus und keine partielle Inner-Konvergenz.

\textbf{Warum das oft fälschlich als alternierend wahrgenommen wird.} In
der Praxis startet die Inner-Suche pro Subjekt mit dem \(\eta_i^*\) aus
der vorherigen Outer-Iteration (Warmstart). Wenn das Outer nahe am
Optimum ist, ist das Inner schon nach 1 bis 2 Newton-Schritten wieder
konvergent. Das \emph{Verhalten} sieht aus wie ``Outer-Schritt,
Inner-Update, Outer-Schritt, Inner-Update'', aber formal ist es strikt
geschachtelt: das Outer wartet auf vollständige Inner-Konvergenz.

\textbf{Abgrenzung zu echten EM-Verfahren:}

{\def\LTcaptype{none} % do not increment counter
\begin{longtable}[]{@{}
  >{\raggedright\arraybackslash}p{(\linewidth - 6\tabcolsep) * \real{0.2500}}
  >{\raggedright\arraybackslash}p{(\linewidth - 6\tabcolsep) * \real{0.2500}}
  >{\raggedright\arraybackslash}p{(\linewidth - 6\tabcolsep) * \real{0.2500}}
  >{\raggedright\arraybackslash}p{(\linewidth - 6\tabcolsep) * \real{0.2500}}@{}}
\toprule\noalign{}
\begin{minipage}[b]{\linewidth}\raggedright
Verfahren
\end{minipage} & \begin{minipage}[b]{\linewidth}\raggedright
Inner-Anteil
\end{minipage} & \begin{minipage}[b]{\linewidth}\raggedright
Outer-Anteil
\end{minipage} & \begin{minipage}[b]{\linewidth}\raggedright
Ablauf
\end{minipage} \\
\midrule\noalign{}
\endhead
\bottomrule\noalign{}
\endlastfoot
\textbf{FOCE / FOCEI} & Mode-Suche
\(\eta_i^* = \arg\max p(y_i\mid\eta_i)p(\eta_i)\) pro Subjekt &
minimiert approx. marginal \(-2\log L\) über \((\theta,\Omega,\Sigma)\)
& \textbf{geschachtelt}, Inner stets voll konvergent \\
Laplace & wie FOCE, mit voller Hessischer & wie FOCE & geschachtelt \\
\textbf{SAEM} & Monte-Carlo-Sampling aus Posterior
\(p(\eta_i\mid y_i,\theta)\) & M-step über \((\theta,\Omega,\Sigma)\) &
\textbf{echt alternierend}, stochastische Approximation, Burn-in dann
Annealing \\
EM (analytisch) & E-step: Posterior-Erwartung & M-step: Maximierung &
alternierend, für NLME aber nicht tractable \\
\end{longtable}
}

\textbf{Was es bei FOCEI also nicht gibt:}

\begin{itemize}
\tightlist
\item
  Keinen Burn-in mit nachgeschalteter Annealing-Phase. Das ist ein
  SAEM-Konzept (Monolix), nötig wegen MC-Rauschen.
\item
  Keine zwei Phasen mit unterschiedlichem Algorithmusverhalten.
\item
  Keine Stochastik im Optimierungspfad (deterministische
  Bilevel-Optimierung).
\end{itemize}

\textbf{Was FOCEI in der Praxis trotzdem als ``Phasen-Verhalten'' sehen
kann:}

\begin{itemize}
\tightlist
\item
  Anfangs (weit weg vom Optimum) braucht das Inner mehr Newton-Schritte
  pro Subjekt, später (nahe Optimum mit Warmstart) sehr wenige. Das ist
  Effizienz-Effekt, keine Algorithmus-Phase.
\item
  Eine \emph{Inner-Toleranz}, die mit der Outer-Gradientennorm skaliert
  (anfangs locker, später strenger), kann die Effizienz weiter
  verbessern. Das ist eine Implementationsoption, kein konzeptioneller
  Phasenwechsel.
\end{itemize}

In dieser Implementierung also: ein einziger geschachtelter Loop,
Warmstart pro Subjekt, optional adaptive Inner-Toleranz. Kein EM, kein
Burn-in.

\subsubsection{2.5 Verwandte Verfahren und
Architekturentscheidung}\label{verwandte-verfahren-und-architekturentscheidung}

Sobald wir die Frage ``warum nicht alternierend?'' konsequent zu Ende
denken, landen wir bei verschiedenen verwandten Methoden, die die
gleiche Laplace-Approximation nutzen, aber unterschiedlich daraus eine
Schätzgleichung bauen. Lohnt sich ein Überblick, weil die
\emph{Mathematik unterhalb} fast identisch ist und wir sie nur einmal
bauen müssen.

\paragraph{Verfahren im Vergleich}\label{verfahren-im-vergleich}

Sei \(J_i(\theta, \eta_i, \Omega, \Sigma)\) die joint negative log
density (Residual + Prior auf \(\eta\), ohne \(\log|H|\)), \(\eta_i^*\)
ihr Mode, \(H_i = \partial^2 J_i / \partial \eta^2\) am Mode.

{\def\LTcaptype{none} % do not increment counter
\begin{longtable}[]{@{}
  >{\raggedright\arraybackslash}p{(\linewidth - 6\tabcolsep) * \real{0.2500}}
  >{\raggedright\arraybackslash}p{(\linewidth - 6\tabcolsep) * \real{0.2500}}
  >{\raggedright\arraybackslash}p{(\linewidth - 6\tabcolsep) * \real{0.2500}}
  >{\raggedright\arraybackslash}p{(\linewidth - 6\tabcolsep) * \real{0.2500}}@{}}
\toprule\noalign{}
\begin{minipage}[b]{\linewidth}\raggedright
Verfahren
\end{minipage} & \begin{minipage}[b]{\linewidth}\raggedright
Outer-Zielfunktion
\end{minipage} & \begin{minipage}[b]{\linewidth}\raggedright
Optimierungsstrategie
\end{minipage} & \begin{minipage}[b]{\linewidth}\raggedright
Ω-Update
\end{minipage} \\
\midrule\noalign{}
\endhead
\bottomrule\noalign{}
\endlastfoot
\textbf{FOCEI} & \(\sum_i [J_i(\eta_i^*) + \log|H_i|]\) (Laplace-Approx
von \(-2\log L_{\text{marginal}}\)) & nested, Outer mit voller Chain
Rule durch \(\eta^*\) & als Teil des Outers numerisch \\
\textbf{Laplace 2. Ordnung} & wie FOCEI plus Korrekturterme aus echter
Hessischer & nested & wie FOCEI \\
\textbf{Laplace-EM} & ELBO
\(\sum_i \bigl[-\tfrac{1}{2} E_{q_i}[J_i] + \tfrac{1}{2}\log|H_i^{-1}|\bigr]\)
mit \(q_i = \mathcal{N}(\eta_i^*, H_i^{-1})\) & alternierend (E-step:
\(\eta_i^*\), \(H_i\); M-step: \(\theta, \Omega, \Sigma\)) &
\textbf{closed form}:
\(\Omega^{(t+1)} = \tfrac{1}{N}\sum_i (\eta_i^* {\eta_i^*}^\top + H_i^{-1})\) \\
\textbf{ITS} & \(\sum_i J_i\) allein (kein \(\log|H|\)) & alternierend,
Block-Koordinatenabstieg & Stichprobenkovarianz von \(\eta_i^*\) \\
\textbf{SAEM} & stochastische ELBO mit MC-Posterior & alternierend, mit
Annealing & closed form aus MC-Samples \\
\end{longtable}
}

Beobachtung: alle Verfahren teilen denselben Inner-Mode \(\eta_i^*\) und
die gleiche Hessische \(H_i\), sie unterscheiden sich nur in (a) der
skalaren Outer-Zielfunktion und (b) der Strategie zum Optimieren. Die
teure Mathematik (ODE-Integration, Sensitivitäten, Cholesky-Ableitungen)
ist gemeinsam.

\paragraph{Wer macht das heute mit exakten
Ableitungen?}\label{wer-macht-das-heute-mit-exakten-ableitungen}

{\def\LTcaptype{none} % do not increment counter
\begin{longtable}[]{@{}
  >{\raggedright\arraybackslash}p{(\linewidth - 4\tabcolsep) * \real{0.3333}}
  >{\raggedright\arraybackslash}p{(\linewidth - 4\tabcolsep) * \real{0.3333}}
  >{\raggedright\arraybackslash}p{(\linewidth - 4\tabcolsep) * \real{0.3333}}@{}}
\toprule\noalign{}
\begin{minipage}[b]{\linewidth}\raggedright
Software
\end{minipage} & \begin{minipage}[b]{\linewidth}\raggedright
Verfahren
\end{minipage} & \begin{minipage}[b]{\linewidth}\raggedright
Sensitivitäten
\end{minipage} \\
\midrule\noalign{}
\endhead
\bottomrule\noalign{}
\endlastfoot
NONMEM & FO, FOCE, FOCEI, Laplace 2. Ordnung & überwiegend FD \\
Monolix & SAEM & überwiegend FD \\
nlmixr2 & FOCEI, SAEM & gemischt, AD nur in Teilen \\
TMB / glmmTMB & Laplace direkt (nicht EM) & volle AD, \emph{aber}
GLMM/lineare Modelle, kein ODE-RHS \\
Pumas (Julia) & FOCEI, Laplace & volle AD inkl. ODE \\
\textbf{dMod-Plan} & FOCEI primär, Laplace-EM optional & \textbf{volle
exakte Ableitungen via CppODE / funCpp} \\
\end{longtable}
}

Die echte Lücke im R-Ökosystem: \textbf{NLME mit ODE-Modellen und
exakten Ableitungen durch die komplette Pipeline (Modell,
Sensitivitäten, FOCEI-Aggregation, Cholesky-Diff)}. Genau das, was dMods
Infrastruktur ohnehin schon hergibt.

\paragraph{Warum FOCEI und nicht Laplace-EM als primäres
Ziel?}\label{warum-focei-und-nicht-laplace-em-als-primuxe4res-ziel}

\textbf{Pro Laplace-EM:}

\begin{itemize}
\tightlist
\item
  \(\Omega\)-Update in geschlossener Form ist mathematisch elegant und
  reduziert die Outer-Dimension um \(K(K+1)/2\) Parameter.
\item
  Konzeptionell sauber alternierend, einfacher Pseudocode.
\item
  Variations-Lower-Bound monoton wachsend (zumindest die ELBO, nicht
  \(\log L\)).
\end{itemize}

\textbf{Contra Laplace-EM:}

\begin{itemize}
\tightlist
\item
  ELBO und Laplace-Marginal-Likelihood sind \textbf{unterschiedliche
  Funktionen}. Konvergenzpunkt ist nicht der FOCEI-Punkt. Das wirft
  Validierungsprobleme gegen NONMEM-Referenzen auf.
\item
  EM-Monotonie greift nur auf der ELBO, nicht auf \(\log L\). Der
  psychologische Vorteil ``EM konvergiert immer'' trügt bei nichtexakter
  Posterior.
\item
  M-step für \(\theta\) und \(\Sigma\) bleibt eine nichtlineare
  Optimierung mit ODE-Auswertungen, also nicht billiger als der
  FOCEI-Outer.
\item
  Pharmakometrie-Welt kennt FOCEI, Reviewer und Regulierungsbehörden
  auch. Laplace-EM ist exotisch und braucht zusätzliche Begründung.
\end{itemize}

\textbf{Pro FOCEI:}

\begin{itemize}
\tightlist
\item
  Standardziel, validierbar gegen NONMEM/nlmixr2/Monolix-Referenzen.
\item
  Lingua Franca der Pharmakometrie.
\item
  Mathematik ist sauberer in einer einzigen Outer-Schleife mit
  \texttt{trust}.
\end{itemize}

\textbf{Contra FOCEI:}

\begin{itemize}
\tightlist
\item
  Cross-Term
  \(\partial \log|H|/\partial \eta \cdot \mathrm{d}\eta^*/\mathrm{d}\theta\)
  in der Outer-Ableitung ist nicht trivial.
\item
  Kein closed-form \(\Omega\)-Update, \(\Omega\) wird im Outer zusammen
  mit \(\theta, \Sigma\) optimiert.
\end{itemize}

Auf das Wesentliche reduziert: FOCEI ist die
\textbf{standardvalidierbare Wahl}, Laplace-EM die \textbf{strukturell
elegantere}.

\paragraph{Architektur: ein Math-Layer, mehrere
Outer-Strategien}\label{architektur-ein-math-layer-mehrere-outer-strategien}

Die saubere Konsequenz: bauen wir die gemeinsame Mathematik einmal
sauber auf, sind FOCEI und Laplace-EM (und auch ITS als
Burn-in-Initialisierung) verschiedene Outer-Strategien auf demselben
Substrat.

\begin{verbatim}
nlme_components(joint, etaSpec, OmegaSpec)
   ->  list(
         inner_solve(theta, Omega, Sigma)    # liefert eta_i*, H_i pro Subjekt
         joint_at_modes(...)                 # J_i(eta_i*) und Ableitungen
         Hi_logdet(...)                      # log|H_i| und Ableitungen
         Omega_chol_algebra(...)             # Inv, det, dInv/dchol etc.
       )

focei(components)        # nested, Outer = trust auf J + log|H|
laplace_em(components)   # alternierend, closed-form Omega update + Outer auf Rest
its(components)          # alternierend ohne log|H|, fuer Burn-in
\end{verbatim}

Der gemeinsame Math-Layer (\texttt{nlme\_components}) macht den teuren
Teil: ODE-Integration, Sensitivitäten, Cholesky-Algebra, alles mit
exakten Ableitungen durch die existierende dMod/CppODE-Pipeline. Die
Outer-Strategien obendrauf sind dünne Wrapper, die nur die jeweiligen
Schätzgleichungen orchestrieren.

\paragraph{Empfehlung}\label{empfehlung}

\begin{enumerate}
\def\labelenumi{\arabic{enumi}.}
\tightlist
\item
  \textbf{FOCEI als primäres Implementierungsziel}, Validierung gegen
  Benchmark-Modelle aus NONMEM und nlmixr2.
\item
  \textbf{Math-Layer so faktoriert, dass Laplace-EM später als
  Outer-Variante reinkommt}. Kosten gering, sobald die FOCEI-Maschinerie
  steht. ITS als optionales Initialisierungsverfahren ebenfalls aus
  diesem Substrat.
\item
  \textbf{Die echte Story} dieser Implementierung gegenüber dem Stand
  der Technik ist \emph{nicht} FOCEI vs.~Laplace-EM, sondern
  \textbf{NLME mit exakten Ableitungen durch alle Schichten}. Das gilt
  für FOCEI und Laplace-EM gleichermaßen, ist aber die wirklich neue
  Lücke im R-Ökosystem.
\end{enumerate}

\begin{center}\rule{0.5\linewidth}{0.5pt}\end{center}

\subsection{3. Existierende
dMod-Bausteine}\label{existierende-dmod-bausteine}

{\def\LTcaptype{none} % do not increment counter
\begin{longtable}[]{@{}
  >{\raggedright\arraybackslash}p{(\linewidth - 2\tabcolsep) * \real{0.5000}}
  >{\raggedright\arraybackslash}p{(\linewidth - 2\tabcolsep) * \real{0.5000}}@{}}
\toprule\noalign{}
\begin{minipage}[b]{\linewidth}\raggedright
Baustein
\end{minipage} & \begin{minipage}[b]{\linewidth}\raggedright
Status
\end{minipage} \\
\midrule\noalign{}
\endhead
\bottomrule\noalign{}
\endlastfoot
ODE-Lösung mit Sensitivitäten & vorhanden via CppODE / deSolve \\
Bedingungsspezifische Parameter & vorhanden via \texttt{parlist},
PEtab-Multi-Condition \\
Komponierbare Pipeline \texttt{g*x*p} & vorhanden, Kettenregel intern \\
\texttt{normL2} Residualbeitrag inkl. Gradient/Hessian & vorhanden \\
\texttt{trust()} Optimizer mit Bound Constraints & vorhanden \\
Parallelisierung über Bedingungen via \texttt{mclapply} & vorhanden \\
Symbolische Jacobians, AD via CppODE & vorhanden
(\texttt{derivMode\ =\ "dual"/"symbolic"}) \\
Profile Likelihoods & vorhanden \\
\end{longtable}
}

Die mentale Umstellung: in dMod ist heute ``Bedingung'' eine
Datengruppe. In FOCEI wird \textbf{Subjekt = Bedingung}, plus eine
populationsweite Kopplung über \(\Omega\), die im aktuellen
\texttt{normL2}-Pfad nicht vorgesehen ist.

\begin{center}\rule{0.5\linewidth}{0.5pt}\end{center}

\subsection{4. Designidee: joint likelihood plus
FOCEI-Wrapper}\label{designidee-joint-likelihood-plus-focei-wrapper}

\subsubsection{4.1 Joint negative log
posterior}\label{joint-negative-log-posterior}

Der Nutzer schreibt das volle penalized-likelihood-Objektiv für alle
Subjekte zusammen: \[
J(\theta, \eta_{1..N}, \Omega, \Sigma)
\;=\; \sum_i \!\left[\, \sum_j \frac{(y_{ij} - f_{ij})^2}{\Sigma_{ij}} + \log\Sigma_{ij}
\;+\; \eta_i^\top \Omega^{-1} \eta_i + \log|\Omega| \,\right]
\]

In dMod-Komposition (Variante A, der Nutzer baut die joint selbst):

\begin{Shaded}
\begin{Highlighting}[]
\NormalTok{joint }\OtherTok{\textless{}{-}} \FunctionTok{normL2}\NormalTok{(data, g }\SpecialCharTok{*}\NormalTok{ x }\SpecialCharTok{*}\NormalTok{ p) }\SpecialCharTok{+} \FunctionTok{constraintL2}\NormalTok{(}\AttributeTok{mu =}\NormalTok{ etas, }\AttributeTok{Omega =}\NormalTok{ OmegaSpec)}
\end{Highlighting}
\end{Shaded}

\texttt{normL2(...)} liefert den Residualanteil mit beliebigem
\texttt{err}-Modell. Der Prior-Beitrag
\(\sum_i \eta_i^\top \Omega^{-1} \eta_i + N\log|\Omega|\) kommt aus
einer Erweiterung von \texttt{constraintL2}: das heutige
\texttt{constraintL2(mu,\ sigma)} ist skalar mit diagonaler \(\sigma\).
Wir geben ihm einen optionalen \texttt{Omega}-Pfad, der bei
MVN-Cholesky-Spec auf \(\eta^\top\Omega^{-1}\eta + \log|\Omega|\) inkl.
Ableitungen nach \(\eta\) und nach den Cholesky-Einträgen umschaltet.
Skelett (objfn-Vertrag, \texttt{dP}-Chain-Rule) bleibt identisch.

\subsubsection{4.2 FOCEI-Wrapper}\label{focei-wrapper}

Eine Funktion, die das Bilevel-Splitting macht:

\begin{Shaded}
\begin{Highlighting}[]
\NormalTok{focei\_obj }\OtherTok{\textless{}{-}} \FunctionTok{focei}\NormalTok{(}
\NormalTok{  joint,                     }\CommentTok{\# objfn auf (theta, eta\_1..N, Omega, Sigma)}
\NormalTok{  etaSpec,                   }\CommentTok{\# welche Parameter sind eta, pro Subjekt}
\NormalTok{  OmegaSpec,                 }\CommentTok{\# welche Parameter parametrisieren Omega (Cholesky)}
  \AttributeTok{innerControl =} \FunctionTok{list}\NormalTok{(...)   }\CommentTok{\# rtol, maxit, warmstart cache}
\NormalTok{)}
\end{Highlighting}
\end{Shaded}

\texttt{focei\_obj} ist wieder ein \texttt{objfn}, aber jetzt nur über
\((\theta, \Omega, \Sigma)\). Bei jedem Aufruf:

\begin{enumerate}
\def\labelenumi{\arabic{enumi}.}
\tightlist
\item
  Für jedes Subjekt \(i\): \texttt{trust()} auf \texttt{joint} bei
  fixiertem \((\theta, \Omega, \Sigma)\), frei in \(\eta_i\). Liefert
  \(\eta_i^*\) und die innere Hessische \(H_i\) zurück.
\item
  FOCEI-Beitrag:
  \(\mathrm{OFV}_i = J_i(\theta, \eta_i^*, \Omega, \Sigma) + \log|H_i|\).
\item
  Summe über \(i\).
\item
  Gradient/Hessian nach \((\theta, \Omega, \Sigma)\) aus der vorhandenen
  \texttt{joint}-Maschinerie plus dem \(\log|H_i|\)-Term.
\end{enumerate}

\subsubsection{4.3 Drei wichtige
Beobachtungen}\label{drei-wichtige-beobachtungen}

\textbf{(a) Inner und Outer sauber separierbar via Envelope Theorem.}
Bei \(\eta_i^*\) ist \(\partial J / \partial \eta = 0\). Damit
verschwindet \(\mathrm{d}\eta_i^*/\mathrm{d}\theta\) aus dem Gradienten
von \(J_i(\theta, \eta_i^*, \dots)\) nach \(\theta\). Wir brauchen nur
die direkten partiellen Ableitungen, die \texttt{joint} ohnehin liefert.
Für die Outer-Likelihood ohne \(\log|H_i|\) ist der Outer-Gradient damit
\emph{trivial}: \texttt{joint} an \((\theta, \eta_i^*, \Omega, \Sigma)\)
auswerten und die Komponenten zu \((\theta, \Omega, \Sigma)\) ablesen.

\textbf{(b) \texttt{trust()} liefert \(H_i\) frei Haus.} Damit ist
\(\log|H_i|\) in der OFV ohne Extraarbeit auswertbar. Schwieriger: die
\emph{Ableitung} von \(\log|H_i|\) nach \((\theta, \Omega, \Sigma)\).
Drei Optionen:

\begin{itemize}
\tightlist
\item
  \(\log|H_i|\) in der Outer-Ableitung schlicht ignorieren.
  Outer-Schritte werden mit dem dominanten Teil gewählt. Konvergenz
  langsamer, aber pragmatisch (analog nlmixr2).
\item
  Numerisch differenzieren in der Outer-Schleife.
  \(\dim(\theta) + \dim(\Omega) + \dim(\Sigma)\) extra Inner-Solves pro
  Outer-Schritt, mit Warmstart bezahlbar.
\item
  Vollanalytisch: bräuchte
  \(\partial^3 f / \partial\eta^2 \partial\theta\). Nicht praktikabel.
\end{itemize}

Siehe Abschnitt 5 (funCpp-Codegen) für die saubere Variante:
kompilierter symbolischer OFV-Knoten mit Gradient durch Chain Rule, der
\(\log|H_i|\) konsistent mitdifferenziert.

\textbf{(c) \(H_i\) in Gauss-Newton-Form genügt.} \[
H_i \;\approx\; G_i^\top \Sigma_i^{-1} G_i + \Omega^{-1} + \text{(Interaktion)}, \qquad G_i = \partial f / \partial \eta_i.
\] Erste-Ableitungs-Information, die das
\texttt{prdframe}-\texttt{deriv}-Attribut bereits liefert. Wir brauchen
\texttt{\_s2} / \texttt{\_sdcv}-Files (zweite Ableitungen aus CppODE)
\textbf{nicht} für FOCEI selbst, nur falls wir später echtes Laplace
machen.

\subsubsection{4.4 Was der Wrapper konkret wissen
muss}\label{was-der-wrapper-konkret-wissen-muss}

\begin{itemize}
\tightlist
\item
  \textbf{\(\eta\)}: pro Subjekt \(K\)-Vektor. Mapping aus
  \texttt{parlist}-Struktur, da Subjekte = Conditions.
\item
  \textbf{\(\Omega\)}: \(K \times K\), symmetric PD, \(K(K+1)/2\) freie
  Parameter. Cholesky \(\Omega = LL^\top\), \(L\) untere Dreiecks,
  \(L_{kk} = \exp(\omega_k)\), \(L_{kl}\) frei für \(k > l\).
\item
  \textbf{\(\Sigma\)}: im \texttt{err}-Modell, das \texttt{normL2} schon
  kennt. In der Outer-Sicht von \(\eta\) und \(\Omega\) trennbar via
  Namens-Partition.
\item
  \textbf{\(\theta\)}: der Rest.
\end{itemize}

Die Aufteilung ist primär eine Namens-Partition über den
Gesamtparametervektor, ähnlich wie heute Bedingungs-Parameter zugewiesen
werden. Mögliche Hilfsfunktion: \texttt{nlmePartition()}.

\begin{center}\rule{0.5\linewidth}{0.5pt}\end{center}

\subsection{5. Codegen via funCpp und 3.
Ableitung}\label{codegen-via-funcpp-und-3.-ableitung}

\subsubsection{5.1 Idee}\label{idee}

Die OFV-Aggregation pro Subjekt \[
\mathrm{OFV}_i = \sum_j \frac{r_{ij}^2}{\Sigma_{ij}} + \log\Sigma_{ij}
\;+\; \eta_i^{*\top} \Omega^{-1} \eta_i^*
\;+\; \log|\Omega|
\;+\; \log|H_i(G_i, \Sigma, \Omega)|
\] ist ein geschlossener symbolischer Ausdruck in den Inputs

\begin{itemize}
\tightlist
\item
  \(f_{ij}\), \(G_i = \partial f / \partial \eta\) (numerischer
  ODE/Sensi-Output zur Auswertungszeit),
\item
  \(\Omega\)-Cholesky (Outer-Param),
\item
  \(\Sigma\)-Params (Outer-Param, beliebige Form),
\item
  \(\eta_i^*\) (aus Inner),
\item
  \(y_{ij}\) (Daten).
\end{itemize}

Genau dafür ist \texttt{funCpp} gemacht. Gleiches Pattern wie heute
\texttt{Y} (Observable) und \texttt{Pexpl} (Parametertrafo):
symbolischer Ausdruck wird mit Ableitungen kompiliert und sitzt als
Knoten in der dMod-Pipeline. Der FOCEI-Wrapper wird dann nur ein
weiterer kompilierter Endknoten.

\subsubsection{5.2 Welche
Ableitungsordnung?}\label{welche-ableitungsordnung}

Aufschlüsselung des Outer-Gradienten
\(\mathrm{d}\mathrm{OFV}/\mathrm{d}\theta\) via Chain Rule durch die
OFV-Inputs:

\begin{itemize}
\tightlist
\item
  Direkter Pfad: braucht \(\partial f/\partial\theta\) (1. Ordnung) und
  \(\partial G/\partial\theta = \partial^2 f / \partial\eta\partial\theta\)
  (gemischte 2. Ordnung).
\item
  Impliziter Pfad via \(\eta_i^*(\theta)\):
  \(\mathrm{d}\eta_i^*/\mathrm{d}\theta = -H_i^{-1} \partial^2 J/\partial\eta\partial\theta\),
  ebenfalls 2. Ordnung gemischt. Envelope eliminiert den Beitrag im
  \(J\)-Anteil, NICHT im \(\log|H|\)-Anteil.
\end{itemize}

\textbf{Outer-Gradient: maximal 2. Ordnung gemischt.} Das liefert CppODE
heute schon via \texttt{\_s2}.

Outer-Hessian: - \(\mathrm{d}^2\mathrm{OFV}/\mathrm{d}\theta^2\) enthält
\(\partial^3 f / \partial\eta\partial\theta^2\), also echte 3. Ordnung.

\textbf{Outer-Hessian exakt: braucht 3. Ordnung.} Trust-Region kommt mit
Gauss-Newton-Approx der Outer-Hessischen üblicherweise gut zurecht
(insb. bei exaktem Gradient).

\subsubsection{5.3 Empfehlung}\label{empfehlung-1}

Stufe 1: exakter Gradient, Gauss-Newton-approximierter Outer-Hessian.
Reicht funCpp mit bestehender 2. Ordnung. Kein Codegen-Aufbohrung.

Stufe 2 (optional, später): falls Konvergenz zickt, funCpp auf 3.
Ordnung erweitern (sowohl symbolisch als auch in der Dual-AD-Variante,
da Forward-Mode-AD kompositionell auf höhere Ordnungen geht).
Eigenständiges Subprojekt im CppODE-Repo.

\begin{center}\rule{0.5\linewidth}{0.5pt}\end{center}

\subsection{6. Verbleibende offene
Designfragen}\label{verbleibende-offene-designfragen}

\begin{enumerate}
\def\labelenumi{\arabic{enumi}.}
\item
  \textbf{Inner-Toleranz adaptiv?} Je näher der Outer am Optimum, desto
  enger muss das Inner konvergiert sein, sonst rauscht der
  Outer-Gradient. Standardrezept (Toleranzfaktor x Outer-Gradientennorm)
  oder konfigurierbar?
\item
  \textbf{Warmstart-Speicher}: Closure-Variable im FOCEI-Wrapper hält
  \(\eta_i^*\) pro Subjekt (bestätigt). Frage: was bei Outer-Reset
  (Multistart, neuer Datensatz)? Reset-Hook im Wrapper?
\item
  \textbf{Parallelisierung über Subjekte}: \texttt{mclapply} über
  Inner-Solves ist trivial, aber die Reihenfolge der Side-Effekte
  (Warmstart-Cache schreiben) muss ggf. ein deterministischer Reduce
  werden.
\item
  \textbf{API-Konvention für \texttt{etaSpec} / \texttt{OmegaSpec}}:
  konkretes Format. Idee:
  \texttt{etaSpec\ =\ list(eta\_kel\ =\ c("subj1",\ "subj2",\ ...),\ eta\_V\ =\ c(...))},
  oder über die existierende \texttt{parlist}-Struktur abgeleitet. Klärt
  sich erst in Phase 1 der Implementation.
\end{enumerate}

\end{document}
